\chapter{Cơ sở lý thuyết và khác biệt khi chuyển đổi}
\section{Vòng lặp huấn luyện và logits}
PyTorch biểu diễn chu trình \texttt{zero\_grad}, forward, loss, backward và optimizer step. Phân loại dùng logits với \texttt{CrossEntropyLoss}; không áp softmax trước loss. Parameter-count parity không tự chứng minh output hoặc metric parity.

Mô hình phân loại xuất logits $z\in\mathbb{R}^{B\times K}$ và class-index target. So sánh với Keras cần ghi nhãn, initialization, reduction và preprocessing; cùng số tham số chưa đủ để khẳng định cùng phép tính.

\section{NHWC sang NCHW}
Keras thường dùng $(N,H,W,C)$, PyTorch convolution nhận $(N,C,H,W)$:
\[
x_{\mathrm{NCHW}}=\operatorname{permute}(x_{\mathrm{NHWC}},0,3,1,2).
\]
WIKI áp dụng phép xoay legacy \texttt{np.rot90(k=3, axes=(1,2))} trên NHWC trước khi đổi layout. Xoay ảnh và hoán vị trục là hai phép toán khác nhau.

\section{CNN-LSTM}
Thay \texttt{TimeDistributed(CNN)} bằng reshape $(B,T,C,H,W)\to(BT,C,H,W)$, chạy CNN, rồi khôi phục $(B,T,F)$ trước LSTM. Tương đương shape không chứng minh tương đương metric.

Phải giữ nguyên thứ tự frame khi gộp và tách chiều. CNN xử lý từng frame, LSTM nhận lại chuỗi có thứ tự. Parameter count CNN-LSTM khác nhau giữa reports T0206/T0207; cần đối chiếu source snapshot và cách đếm bias/buffer trước khi chốt.

\section{SegNet và upsampling}
Đối chiếu \texttt{UpSampling2D} với \texttt{nn.Upsample}; PyTorch dùng logits phân lớp từng pixel và mask class-index cho cross-entropy. Báo riêng total/trainable/non-trainable theo quy ước Keras.

Với batch ảnh, logits thường có shape $(B,K,H,W)$, còn class-index mask có shape $(B,H,W)$. Pixel accuracy và mean IoU không phải một metric; giữ đúng tên/chỉ số có bằng chứng.

\section{FCN score-map và chuẩn hóa}
Giữ chuẩn hóa legacy $x/127.5-1.0$ khi đúng đầu vào. Score-map regression và metric notebook tên accuracy không tự trở thành IoU/mAP.

Nếu chuyển score grid thành tọa độ, ghi rõ stride và phép ánh xạ. Heatmap hoặc output từ model chưa huấn luyện không chứng minh detection quality.

\section{Chuẩn hóa và giao thức đánh giá}
Nêu preprocessing theo dataset. Nếu validation/public-test được tái sử dụng, kết quả không đại diện independent final holdout.

\section{Các yếu tố ảnh hưởng đến parity}
Kiểm tra riêng initializer, padding/pooling, Dropout ở train/eval, BatchNorm statistics/epsilon/momentum, layout, thứ tự class labels, loss reduction, optimizer và seed. Parameter count chỉ kiểm tra cấu trúc một phần.